\begin{table*}[t]
\centering
\footnotesize
\setlength{\tabcolsep}{5pt}
\begin{tabular}{l r r l l l}
\toprule
Task & Items & Cells & Original window & Repaired window & Change \\
\midrule
convert vs.\ preference & 2 & 24 & 0.000 [0.000, 0.138] & 0.083 [0.023, 0.258] & $-$0.083 [$-0.258$, $0.067$] \\
multi-round dilution & 5 & 60 & 0.167 [0.093, 0.280] & 0.200 [0.118, 0.318] & $-$0.033 [$-0.172$, $0.107$] \\
option pool shuffle & 3 & 36 & 0.611 [0.449, 0.752] & 0.306 [0.180, 0.469] & 0.306 [$0.075$, $0.494$] \\
\midrule
\textbf{Pooled} & 10 & 120 & 0.267 [0.196, 0.352] & 0.208 [0.145, 0.289] & 0.058 [$-0.049$, $0.165$] \\
\bottomrule
\end{tabular}
\caption{The re-windowing repair, 10 items each measured in both conditions by all twelve configurations (120 paired cells). The repaired window is cut to the same length as the original and centred on the validating quote, so the only difference is whether the evidence is inside it. Registered before the experiment ran: the repaired rate would fall below 0.161. It is 0.208 [0.145, 0.289], above that threshold, so the prediction fails. The pooled change is 0.058 [$-0.049$, $0.165$] and contains zero. Change is the original rate minus the repaired one, so a positive value is a fall in wobble and a negative value a rise; 2 of the 3 tasks are negative. One task of three moves; the per-task split is descriptive and was not the registered test.}
\label{tab:repair}
\end{table*}
